\section*{Capítulo \ref{ch:g10:freshwater-resources} --- Los recursos de agua dulce}

\begin{solution}{exo:g10:freshwater-resources:1}
Una \omterm{def:g10:freshwater-resources:stock}{reserva} es la cantidad de agua contenida en un depósito en un momento
dado, en \unit{km^3} (o en kilogramos). Un \omterm{def:g10:freshwater-resources:stock}{flujo} es la cantidad que pasa de un
depósito a otro por unidad de tiempo, en \unit{km^3} por año. El \omterm{def:g10:freshwater-resources:residence-time}{tiempo de
residencia} es el tiempo medio que el agua pasa en un depósito, en días o en
años; en estado estacionario es la \omterm{def:g10:freshwater-resources:stock}{reserva} dividida entre el \omterm{def:g10:freshwater-resources:stock}{flujo}.
\end{solution}

\begin{solution}{exo:g10:freshwater-resources:2}
Alrededor del \qty{2.5}{\percent} de toda el agua es dulce. Del agua dulce,
el \qty{68.7}{\percent} está helado en casquetes de hielo, glaciares y nieves
perpetuas, y el \qty{30.1}{\percent} está bajo tierra. Los ríos son los que
menos contienen, \num{2120}~\unit{km^3}.
\end{solution}

\begin{solution}{exo:g10:freshwater-resources:3}
El \omterm{def:g10:freshwater-resources:renewable}{agua dulce renovable} es el agua dulce que la precipitación renueva cada
año y que atraviesa una región como escorrentía y como agua subterránea: un
volumen por año, un \omterm{def:g10:freshwater-resources:stock}{flujo}. El \omterm{def:g10:freshwater-resources:fossil-groundwater}{agua subterránea fósil} es agua que entró en un
acuífero hace mucho tiempo y que la recarga actual no renueva: es una
\omterm{def:g10:freshwater-resources:stock}{reserva} que solo puede ir mermando.
\end{solution}

\begin{solution}{exo:g10:freshwater-resources:4}
$91\,000/2120 \approx 43$ veces más.
\end{solution}

\begin{solution}{exo:g10:freshwater-resources:5}
La atmósfera es la que renueva su agua más deprisa (unos \num{9.5} días), y
el océano, el más lento de los depósitos calculados (unos \num{3200} años);
el agua subterránea profunda es más antigua todavía. $3200 \times 365.25/9.5 \approx \num{120000}$ veces más tiempo.
\end{solution}

\begin{solution}{exo:g10:freshwater-resources:6}
Entradas: $413\,000 + 73\,000 = \num{486000}$~\unit{km^3} al año, luego
$T = 12\,700/486\,000 \approx \num{0.0261}$ años $\approx \num{9.5}$ días.
Salidas: $373\,000 + 113\,000 = \num{486000}$~\unit{km^3}, lo mismo: el aire
está en estado estacionario, y el \omterm{def:g10:freshwater-resources:residence-time}{tiempo de residencia} es el mismo.
\end{solution}

\begin{solution}{exo:g10:freshwater-resources:7}
$1\,335\,040/413 \approx \num{3230}$ años. En estado estacionario, el océano
recibe tanto como pierde: $373 + 40 = 413$ miles de~\unit{km^3} al año, así
que el \omterm{def:g10:freshwater-resources:residence-time}{tiempo de residencia} calculado con las entradas es el mismo.
\end{solution}

\begin{solution}{exo:g10:freshwater-resources:8}
Las tierras emergidas reciben \num{113000}~\unit{km^3} al año; \num{73000}
vuelven al aire por evaporación y transpiración y \num{40000} van a parar al
mar: $73\,000 + 40\,000 = \num{113000}$. La parte que vuelve al aire es
$73/113 \approx \qty{65}{\percent}$.
\end{solution}

\begin{solution}{exo:g10:freshwater-resources:9}
$273.2 \times 10^6 \times 1233.5 \approx \qty{3.37e11}{m^3}$, es decir,
\qty{337}{km^3}: casi tanto como toda el agua que el mar de Aral todavía
contenía en 1987.
\end{solution}

\begin{solution}{exo:g10:freshwater-resources:10}
$15.8 \times 0.3048 \approx \qty{4.82}{m}$. A lo largo de los 65 años que van
de 1950 a 2015, $482/65 \approx \qty{7.4}{cm}$ al año.
\end{solution}

\begin{solution}{exo:g10:freshwater-resources:11}
$2/0.01 = 200$ años. La \omterm{def:g10:freshwater-resources:stock}{reserva} baja ahora $0.05 - 0.01 =
\qty{0.04}{km^3}$ al año, así que dura $2/0.04 = 50$ años.
\end{solution}

\begin{solution}{exo:g10:freshwater-resources:12}
1960: $1093/71.1 \approx 15$ años; 2010: $98.1/11.4 \approx 8.6$ años. La
evaporación quita cada año más o menos el mismo espesor de agua de cada
kilómetro cuadrado del lago; al encogerse el lago, su superficie se redujo, y
con ella el volumen total evaporado.
\end{solution}

\begin{solution}{exo:g10:freshwater-resources:13}
(1) La pérdida neta es de $1.2 - 0.2 = \qty{1.0}{m}$ al año sobre
$100 \times 10^6$~\unit{m^2}: tienen que entrar $10^8$~\unit{m^3},
\qty{0.1}{km^3} al año. (2) El aporte baja a \qty{0.05}{km^3}: el lago pierde
$\qty{0.05}{km^3} = 5 \times 10^7$~\unit{m^3}, una capa de
$5 \times 10^7/10^8 = \qty{0.5}{m}$. (3) Al bajar el nivel, las orillas
salen del agua y la superficie del lago se reduce: menos superficie, menos
evaporación, hasta que la evaporación y el aporte vuelven a equilibrarse.
\end{solution}

\begin{solution}{exo:g10:freshwater-resources:14}
El acuífero contiene $100 \times 0.20 = \qty{20}{m}$ de agua, es decir,
\qty{20000}{mm}: a \qty{1}{mm} al año, \num{20000} años. El agua profunda del
acuífero del Sahara tiene entre \num{200000} años y un millón de años, de diez
a cincuenta veces más: se recargó durante climas pasados más húmedos.
\end{solution}

\begin{solution}{exo:g10:freshwater-resources:15}
$337/453\,000 \approx \qty{0.00074}{km}$, unos \qty{0.74}{m} de agua. El
nivel del agua bajó \qty{4.8}{m}, unas seis veces y media más. El agua solo
ocupa los poros de la roca (y no toda se escurre): cuando el nivel freático
baja un metro, el agua liberada es una capa mucho más delgada que un metro,
aquí de unos $0.74/4.8 \approx \num{0.15}$~m.
\end{solution}

\begin{solution}{pb:g10:freshwater-resources:1}
\textbf{1.} $1093/68\,000 \approx \qty{0.016}{km}$, unos \qty{16}{m}.

\textbf{2.} $1093 \times 10^9 \approx \qty{1.09e12}{m^3}$.

\textbf{3.} \qty{10}{g/L} son \qty{10}{kg/m^3}: $1.093 \times 10^{12} \times 10
\approx \qty{1.1e13}{kg}$, unos \num{11}~mil millones de toneladas.

\textbf{4.} $1093/71.1 \approx 15$ años.

\textbf{5.} No sale de él ningún río: su agua solo puede irse como vapor (o
infiltrándose en el suelo), mientras que la sal se queda.

\textbf{6.} $71.1/68\,000 \approx \qty{0.00105}{km}$, unos \qty{1.05}{m}.

\textbf{7.} $9.4/68\,000 \approx \qty{0.00014}{km}$, unos \qty{0.14}{m}.

\textbf{8.} $71.1 - 9.4 = \qty{61.7}{km^3}$ al año.

\textbf{9.} $79 + 37 = \qty{116}{km^3}$: $61.7/116 \approx \qty{53}{\percent}$,
más o menos la mitad.

\textbf{10.} Entradas: $1.0 + 6.2 = \qty{7.2}{km^3}$; salida:
\qty{36.8}{km^3}; variación neta: $7.2 - 36.8 = \qty{-29.6}{km^3}$.

\textbf{11.} $(502.7 - 327.2)/5 = \qty{35.1}{km^3}$ al año: del mismo orden
que los \qty{29.6}{km^3} del balance de 1987 (los años sueltos difieren: la
evaporación fue de \qty{47.9}{km^3} en 1984).

\textbf{12.} $36.8/0.00105 \approx \qty{35000}{km^2}$.

\textbf{13.} $35\,000/68\,000 \approx \qty{52}{\percent}$, más o menos la
mitad.

\textbf{14.} El lago se había encogido: la mitad de la superficie evapora más
o menos la mitad del agua.

\textbf{15.} $345.6/1093 \approx \qty{32}{\percent}$, alrededor de un tercio.

\textbf{16.} $10 \times 1093/345.6 \approx \qty{31.6}{g/L}$.

\textbf{17.} La salinidad medida es menor: parte de la sal había salido del
agua, cristalizada sobre el lecho al descubierto y arrastrada por el viento
en forma de polvo salado.

\textbf{18.} Alrededor de \qty{1.05}{m} al año, el espesor evaporado.

\textbf{19.} $7.2/35\,000 \approx \qty{0.0002}{km}$, unos \qty{0.20}{m}.

\textbf{20.} $1.05 - 0.20 \approx 0.85$ (\num{0.84} con los valores sin
redondear): \textbf{en 1987 el nivel del mar de Aral bajó unos \qty{0.8}{m},
cerca de un metro en un solo año}.
\end{solution}
